\documentclass[twoside,dvipdfm]{article}

\usepackage{htmx}


%<![HTML[
%    <style>
%    <!--
%    h1,h2,h3,h4 { font-family : verdana,helvetica,sans-serif }
%    pre,code    { color: #202080 }
%    body        { padding: 3%  }
%    li          { margin-bottom: 0.5ex }
%    -->
%    </style>
%]]>




\renewenvironment{verbatim}{\small\begin{verbatim}}{\end{verbatim}}


\begin{document}
\title{Babel, integrating languages into Visual Studio}


\author{DAAN LEIJEN
\\
University of Utrecht\strut \\ 
Dept. of Computer Science\strut \\
PO.Box 80.089, 3508 TB Utrecht\strut \\ 
The Netherlands\strut \\
\email{daan@cs.uu.nl}, 
\url{http://www.cs.uu.nl/~daan}
}\maketitle


\hypertarget{Introduction}{\section{Introduction}}\label{Introduction}

This document describes how to integrate programming languages into Visual Studio using 
the {\itshape{}Babel} package. The {\itshape{}Babel package} is available
for both the Visual J++ 6.0 environment and Visual Studio 7.0.
By a language integration we mean: 
Syntax highlighting (including customizable color classes);
Syntax checking, show red squiglies in the source
and show messages in a special task window;
Statement completion; Method tips,  show the type of the method when 
providing parameters; Member selection lists; Quick info, show information
when hovering the mouse over identifiers; Highlight matching braces; find
matching braces; Browsing; Build and deploy management; Project support and
Debugger integration.

There are basically three ways to integrate a language into Visual Studio:
\begin{itemize}
\item{}The Visual Studio Integration Partner (VSIP) program. By becoming a member
of the VSIP program, you get the VSIP SDK which exposes all the core API's of
the Visual Studio environment. Allthough the most flexible and powerful method
of integration, it requires a lot of work to integrate a language. 
\item{}The \texttt{IBabelService} interface. The \texttt{IBabelService} interface 
represents the essential functionality that a language should provide in order to
be integrated into Visual Studio. The {\itshape{}Babel} package provides all
the glue code between this interface and the Visual Studio environment. 
Creating an implementation of the \texttt{IBabelService} is much easier
than using the VSIP SDK and might be the right approach if your language
compiler is easily adapted to support this interface.
\item{}The {\itshape{}Babel} SDK. This SDK implements a default \texttt{IBabelService}
interface in C++ which is customized through a lexical and grammar specification of the language. 
A client only needs to provide these specifications to integrate into Visual Studio.
This is by far the easiest way to create a language service and this document 
describes how to use this approach.
\end{itemize}

The {\itshape{}Babel} package supports (automatic) method tips, 
(automatic) member selection, statement completion,
quick info (automatic info when hovering the
mouse over an identifier), (automatic) highlight matching brace, find matching
brace and block (un)commenting.
We hope to extend the 
{\itshape{}Babel} package soon with browsing, a build system with automatic 
dependency analysis, projects and help integration.


\hypertarget{Getting started}{\section{Getting started}}\label{Getting started}

Figure \href{#fig:overview}{[fig:overview]} shows how the {\itshape{}Babel} 
framework is works. The dynamic link library \texttt{babelpackage.dll} 
is a Visual Studio package that interfaces 
between {\itshape{}Babel} services and Visual Studio. A {\itshape{}Babel} service
is simply a COM component that implements the \texttt{IBabelService} interface
and is registered in a special way. 
This document will explain how to implement such service using the SDK.

\hypertarget{Installation}{\subsection{Installation}}\label{Installation} 
After installing the {\itshape{}Babel} SDK, there is 
a root folder called \texttt{babel\char92{}client}. The root folder contains the following 
sub folders:

\begin{itemize} 
\item{} \texttt{client\char92{}doc}. Contains the documentation about 
{\itshape{}Babel}, FLEX and BISON.

\item{} \texttt{client\char92{}tools}. Contains FLEX and BISON. Both are free GNU 
programs. FLEX generates lexical analysers and BISON generates parsers\footnote{The 
output of both programs is completely free without any licensing issue.}. 

\item{} \texttt{client\char92{}common}. This folder contains the sources of the 
\texttt{common.lib} library. This library implements the standard code 
that implements the \texttt{IBabelService} and interfaces with the generated
lexer and parser.

\item{} \texttt{client\char92{}template}. This is the simplest possible language 
service. This project is used as a template for implementing a language service.

\item{} \texttt{client\char92{}myc} and \texttt{client\char92{}forth}.
These folders 
contain examples of a language integrated into Visual Studio using the 
{\itshape{}Babel} framework. 
\end{itemize}

\hypertarget{Conventions}{\subsection{Conventions}}\label{Conventions}

The interface to the {\itshape{}Babel} language package is specified in 
\texttt{client\char92{}common\char92{}babelservice.idl}. 

The code in \texttt{client\char92{}common} uses the following conventions:

\begin{itemize}
\item{} Types and constants start with a {\bf{}C}apital letter; \texttt{class 
StdService, enum Color \char123{} ColorBlack \char125{}}.

\item{} Variables start with a {\bf{}l}ower case letter; \texttt{token, 
g\_service}. Global variables are prefixed with \texttt{g\_} and member 
variable are prefixed with \texttt{m\_}. 

\item{} Argument use is marked as \texttt{in}, \texttt{out} or 
\texttt{inout}; \texttt{void increment( inout int* i );}. All three 
identifiers are empty \#defines.

\item{} Methods are marked as \texttt{final}, \texttt{virtual} or 
\texttt{override}. \texttt{final} is used for a non-virtual method. 
\texttt{override} is used whenever a virtual method is redefined. Both are 
empty \#defines. 
\end{itemize}

\hypertarget{Adding a language}{\section{Adding a language}}\label{Adding a language}

Adding a language using the {\itshape{}Babel} SDK is very much an evolutionary
process, adding new features one at a time:
\begin{enumerate}
\item{} Set the language information. This step doesn't add any functionality
but enables the integration into Visual Studio.
\item{} Syntax highlighting. This is enabled through a lexical specification.
\item{} Syntax checking. This is enabled trhough a grammar specification.
\item{} Custom colors. Some languages need extended color sets.
\item{} Automatic brace matching. 
\item{} Block commenting.
\item{} Statement completion.
\item{} Member completion.
\item{} Quick info.
\item{} Method tips.
\end{enumerate}

The following sections will elaborate on each step. 
As a running example I will use the {\itshape{}MyC} 
language. {\itshape{}MyC} is an extremely simple C-like language which
only uses integers. 

\hypertarget{Set the language information}{\section{Set the language information}}\label{Set the language information}

This step installs all the basic machinery needed by a language service.

\hypertarget{Create a new project}{\subsection{Create a new project}}\label{Create a new project} 

First a new project for your language is needed. 
Luckily, there is a template project with all the boilerplate code and 
correct settings provided by the SDK. 
Copy the subfolder \texttt{client\char92{}template} to a new subfolder with
the name of your language, for example \texttt{client\char92{}myc}
Click now on 
the \texttt{bservice.dsw} workspace file in your new
subfolder to startup the Visual Studio environment. 

\hypertarget{Language information}{\subsection{Language information}}\label{Language information} 

Open the \texttt{service.cpp} file (in 
{\bf{}source files}). At the top of this file three values need to be 
supplied that are used during the registration of the package:

\begin{itemize}    
\item{} \texttt{g\_languageCLSID} is the CLSID of this specific language. 
The \texttt{guidgen} utility (which can be run from
the \texttt{start-run} menu) is used to generate a fresh CLSID.    
\item{} 
\texttt{g\_languageName} is the human readable name of the language, in our 
case \texttt{"MyC"}.
\item{} \texttt{g\_languageExtensions} are the extensions of the source files
ending with a \texttt{NULL}, for MyC this is simply \texttt{\char123{}".myc",NULL\char125{}}\end{itemize}

For the {\itshape{}MyC} language, the definitions look like:

\begin{alltt}
  const char*  g\_languageName              = "MyC";
  const char*  g\_languageFileExtensions[]  = \char123{} ".myc", NULL \char125{};
  const CLSID  g\_languageCLSID = 
     \char123{}0x54D21404,0x5E7F,0x11D3,\char123{}0xB4,0x07,0x00,0x60,0x08,0xD1,0xBF,0x8C\char125{}\char125{};
\end{alltt}

From this moment on it is safe to build the service, allthough 
no functionality is yet implemented.

\hypertarget{Language properties}{\subsection{Language properties}}\label{Language properties}

The language properties are defined in the \texttt{g\_languageProperties}
variable. This is an array of \texttt{LanguageProperty} structures.
Each property consists of a name and a numeric value. The properties
will be put in the registry on registration and determine the behaviour
of certain features in the language service. For example, automatic
brace matching is only enabled if the property \texttt{"MatchBraces"}
is set to one. When no features are implemented yet, the set of
properties should probably look like:
\begin{alltt}
const LanguageProperty g\_languageProperties[] =
\char123{}
  \char123{} "RequestStockColors", 1     \char125{},
  \char123{} "ShowCompletion",     1     \char125{},

  \char123{} "CodeSense",          0     \char125{},
  \char123{} "MaxErrorMessages",   5     \char125{},
  \char123{} "QuickInfo",          0     \char125{},
  \char123{} "MatchBraces",        0     \char125{},

  \char123{} NULL, 0 \char125{}
\char125{};
\end{alltt}

\hypertarget{Syntax highlighting}{\section{Syntax highlighting}}\label{Syntax highlighting}

Syntax highlighting is enabled through a lexical specification. The 
specification consists of regular expressions that recognise the tokens
of the language. These tokens are than mapped to color classes which
are used by Visual Studio to highlight the source. 
Implementing syntax highlighting is done in three steps:
\begin{itemize}
\item{}Define the tokens of the language.
\item{}Define a mapping from tokens to color classes.
\item{}Define the lexical specication.
\end{itemize}

\hypertarget{Define the tokens of the language}{\subsection{Define the tokens of the language}}\label{Define the tokens of the language} 

Open \texttt{parser.y} (in {\bf{}source files}). Language tokens 
are defined using the \texttt{\%token} definitions of BISON. The 
{\itshape{}MyC} language contains the following definitions:

\begin{alltt}
  \%token IDENTIFIER NUMBER 
  \%token KWIF KWELSE KWWHILE KWFOR KWCONTINUE KWBREAK KWRETURN
  \%token KWEXTERN KWSTATIC KWAUTO KWINT KWVOID
  \%token ',' ';' '(' ')' '\char123{}' '\char125{}' '=' 
  \%token '+' '-' '*' '/' '!' '\&' '|' '\char94{}'
  \%token EQ NEQ GT GTE LT LTE AMPAMP BARBAR 
  \%token LEX\_WHITE LEX\_COMMENT LEX\_LINE\_COMMENT
\end{alltt}

The service should now be build. A 
file called \texttt{parser.cpp.h} is automatically generated by 
{\itshape{}bison} that contains the {\itshape{}C} definitions of the tokens. 

\hypertarget{Define a mapping from tokens to color classes}{\subsection{Define a mapping from tokens to color classes}}\label{Define a mapping from tokens to color classes} 

The tokens 
in the language should be mapped to a specific color class. For example, the 
token \texttt{KWIF} should map to the \texttt{ClassKeyword} color class, a 
\texttt{NUMBER} token to the \texttt{ClassNumber} color class, etc. 

In \texttt{service.cpp}, this information is provided by the 
\texttt{Service::getTokenInfo} method. This method returns a pointer to a 
static \texttt{TokenInfo} array. There is an entry for each defined token. 
Each \texttt{TokenInfo} structure contains the token, the color class for 
this token, a description, a character class and sometimes a trigger class.
The character classes and trigger classes will be explained later in
more detail.
Part of the table for the MyC language looks like:

\begin{alltt}
  static TokenInfo tokenInfoTable[] =
  \char123{}
    //TODO: add your own token information here
    \char123{} IDENTIFIER, ClassIdentifier,  "identifier '\%s'", CharIdentifier \char125{},
    \char123{} NUMBER,     ClassNumber,      "number ('\%s')"  , CharLiteral    \char125{},
    \char123{} KWIF,       ClassKeyword,     "if"             , CharKeyword \char125{},
    \char123{} KWELSE,     ClassKeyword,     "else"           , CharKeyword \char125{},
    \char123{} KWWHILE,    ClassKeyword,     "while"          , CharKeyword \char125{}, 
  
    \char123{} LEX\_WHITE,  ClassText,        "white space"    , CharWhiteSpace \char125{},
    \char123{} LEX\_LINE\_COMMENT, ClassComment, "comment"      , CharLineComment \char125{},
    \char123{} LEX\_COMMENT,ClassComment,     "comment"        , CharComment \char125{},

    //always end with the 'TokenEnd' token.
    \char123{} TokenEnd,     ClassText,      "<unknown>" \char125{}
  \char125{};
\end{alltt}

The color classes of MyC are only the default ones. 
However {\itshape{}Babel} allows for custom 
color classes. The 
default color classes are defined in \texttt{client\char92{}common\char92{}babelservice.idl}:

\begin{alltt}
  enum DefaultColorClass
  \char123{}
    ClassText,
    ClassKeyword,
    ClassComment,
    ClassIdentifier,
    ClassString,
    ClassNumber
  \char125{};
\end{alltt}

\hypertarget{The lexical specification}{\subsection{The lexical specification}}\label{The lexical specification} 

Finally, the lexical specification of the language 
can be added. This specification is used directly by the colorizer to provide 
syntax high-lighting. The same specification is also used later by the 
syntax-checker to provide tokens for the grammar. 

Open the file 
\texttt{lexer.lex} in {\bf{}source files} and add the LEX 
rules for your language. The document \url{flex.html} gives more 
information on how to specify lexical rules. Most languages already have a LEX 
specification. If you have one, you can simply copy the rules into the 
\texttt{lexer.lex} file. However, to meet the requirements of the colorizer, 
there are some constraints on the form of the LEX rules. 

To explain the constraints, we will briefly look into the implementation of the colorizer.
The colorizer 
automatically saves the state of the lexer at every line of text. Whenever the 
line colors could change, the lexer is restarted at this state to colorize this 
line again. Each time that a token is returned by the lexer, the colorizer will 
lookup the color class in the table provided by the \texttt{getTokenInfo} 
method. All the characters scanned for this token are colored with this color 
class. In order to make this work, the colorizer assumes two things from the 
lexical specification:

\begin{itemize}
\item{} The lexer returns a token at the end of a line.
\item{} No scanned characters are dropped or added during scanning.
\end{itemize}

The first constraint assures that the lexer is in a known state at the end of 
the line. The second constraint is needed to determine how many characters are 
to be colored. The first constraint can be easily satisfied by adding a rule at 
each lexer state that returns a token on seeing a new-line. For example, many 
lexical specifications contain the following rule for discarding white space:

\begin{alltt}
  [ \char92{}t\char92{}r\char92{}n]+    \char123{} /* ignore */ \char125{}
\end{alltt}

Since this rule does not directly return a token on a new-line, it should be redefined to:

\begin{alltt}
  [ \char92{}t\char92{}r]+    \char123{} /* ignore */ \char125{}
  \char92{}n          \char123{} /* ignore */ \char125{}
\end{alltt}

However, now it doesn't satisfy the second constraint since it forgets
about the scanned characters. The final and correct definition returns
a white space token:

\begin{alltt}
  [ \char92{}t\char92{}r]+    \char123{} return LEX\_WHITE; \char125{}
  \char92{}n          \char123{} return LEX\_WHITE; \char125{}
\end{alltt}

Since the token mapping specifies the character class \texttt{CharWhiteSpace}
for \texttt{LEX\_WHITE}
only the  colorizer will know about the white space, the parser ignores all white 
space tokens. 
(You can change this behaviour by overriding the \texttt{Service::isGrammarToken}
method).

The second constraint can be enforced in general by always adding an action 
to a lexical rule (so we don't lose characters) and to allow only three ways to 
end that action:

\begin{enumerate}
\item{} \texttt{return {\itshape{}Token};}. "I have scanned a token and return its 
name". This is the normal way of ending an action.

\item{} \texttt{yymore();}. "I have scanned something but don't want to return 
a token yet; continue scanning". \texttt{yymore()} is used here to avoid 
losing characters in \texttt{yytext} (the currently scanned text) 
while continuing scanning. It is used 
in situations like comments where we only want to return the comment token at 
the end of a comment or line.

\item{} \texttt{yyless(0);}. "I have scanned something and changed my state, 
scan this input again". Here, \texttt{yyless(0)} is used to prevent 
duplication of characters. \end{enumerate}

During lexing, there is a global variable \texttt{g\_service} available. This 
variable points to the current language service. The complete interface of a 
language service can be found in \texttt{client\char92{}common\char92{}babelservice.idl} 
and \texttt{client\char92{}common\char92{}stdservice.h}. Important functions during 
lexical analysis are:

\begin{itemize}

\item{} \texttt{void setLexState( in LexState state);}. \strut \\
Sets the lexer state. Should be used instead of the FLEX \texttt{BEGIN} macro.

\item{} \texttt{void lexicalError( in Severity sev, in const char* msg );}. \strut \\
Used to send a lexical error. These messages show  up at run-time in the task 
window. Note that this method has only effect during syntax checking and will 
not work yet during colorizing. An important kind of error is the `hint' which
is used to set for example \texttt{TODO} messages in the task list. A hint is
given by setting the severity to \texttt{SevHint}. The error message should
start with the hint token, not including the comment start. Take a look at
MyC for an example.

\item{} \texttt{void enterComment(in LexState commentState = 0);}. \strut \\
If the \texttt{commentState} argument is given, the lexer is put into that 
state and the current state is saved. Each time this method is called, the 
comment nesting level is incremented.

\item{} \texttt{void leaveComment();}. \strut \\
Each time this method is called, the comment nesting level is decremented. 
Whenever it reaches zero, the lexer is put into the state where 
\texttt{enterComment} was called the first time. 
\end{itemize}

After adding the lexical rules, the colorizer will be working. Open some source 
files of your language in the Visual Studio environment and enjoy them in multiple
colors!

\hypertarget{Syntax checking}{\section{Syntax checking}}\label{Syntax checking}

Open the file 
\texttt{parser.y} again and add the grammar rules for the language. The document
\url{bison.html} contains more information about specifying 
grammar rules. Most languages already have a YACC compatible grammar. You can 
simply copy these rules into \texttt{parser.y}. Unlike the lexer, there are 
no restrictions on the grammar rules. Here are some 
(simplified) rules for {\itshape{}MyC} expressions:

\begin{alltt}
Expr: MulExpr AddOp Expr
    | MulExpr
    ;

AddOp
    : '+' | '-'
    ;


MulExpr
    : PreExpr MulOp MulExpr
    | PreExpr 
    ;

MulOp 
    : '*' | '/'
    ;


Factor
    : IDENTIFIER
    | NUMBER
    | '(' Expr ')'
    ;     
\end{alltt}

\hypertarget{Properties}{\subsection{Properties}}\label{Properties}
After having added the language grammar rules, you should set the
language property \texttt{CodeSense} to one (in \texttt{service.cpp})
to enable syntax checking. The property \texttt{CodeSenseDelay} specifies
the number of milliseconds delay before errors are displayed. 
In practice a value between 1000 or 2000 milliseconds works well, the
default value is 1500.
The property \texttt{CodeSenseFastOnChangeLine} should be set to one
if errors should be displayed immediately when the user leaves a changed
line. This property is enabled by default\footnote{This feature is not 
supported in the Visual J++ 6.0 environment}.
The property \texttt{MaxErrorMessages} gives the maximum number of
displayed error messages (per source file). The default value is 5.

\hypertarget{Adding error productions}{\subsection{Adding error productions}}\label{Adding error productions}
Now, error productions should be added in order to provide useful feedback 
to the user. Error productions should be added in such a way 
that they both provide specific information about the error and, most important, 
that they allow parsing to {\itshape{}recover}. The global language service 
(\texttt{g\_service}) can be used to send error messages about errors that 
are detected during parsing: 

\begin{alltt}
void syntaxError(  in const char*     construct, 
                   in const Location* loc = NULL );

void expectError(  in const char*     construct,
                   in const char*     expecting,
                   in const Location* loc = NULL );

void errorMessage( in Severity        severity,
                   in const char*     message,
                   in const Location* loc = NULL );
\end{alltt}

\texttt{syntaxError} is used for general errors during parsing. For example, 
you probably want to add this 'catch all' error production to your grammar:

\begin{alltt}
SourceFile: Program
          | error     \char123{} g\_service->syntaxError( "program" ); \char125{}
          ;
\end{alltt}

\texttt{expectError} is used whenever a certain closing token is expected. 
For example, in the {\itshape{}MyC} binding there is an error production for 
catching situations where the closing parenthesis of an expression is absent:

\begin{alltt}
ParenExpr
    : '(' Expr ')'    
    | '(' Expr error  
              \char123{} g\_service->expectError( "unmatched parenthesis", ")" ); \char125{}
    ;
\end{alltt}

Sometimes, the location of the error token is not determined by the current 
point in the input. In that case, the location can be given using the 
\texttt{\$} notation of BISON which specifies which token determines 
the position.

\begin{alltt}
ImportStatement
  : IMPORT QualifiedName SemiColons
  | IMPORT QualifiedName '.' '*' SemiColons
  | IMPORT error SemiColons        
       \char123{} g\_service->syntaxError( "import  declaration", \&\$2 ); \char125{}
  ;
\end{alltt}

Putting error productions in the right spot is quite hard, many times error 
productions add a lot of shift-reduce (not so bad) or reduce-reduce (very bad!) 
conflicts to the grammar. Try to aim for adding error productions for `common' 
errors, together with some `catch-all' productions. The MyC sample contains
lots of common error productions.

\hypertarget{Custom color classes}{\section{Custom color classes}}\label{Custom color classes}

For most languages, the default color classes are good enough but some
languages have more token categories. For example, {\itshape{}Haskell} has special
color classes for constructors, literate text and non-standard keywords.
I will use the \texttt{client\char92{}haskell} language as the example for adding
custom color classes.

First of all, we need to specify that we are going to use custom color classes.
Set the language property \texttt{"RequestStockColors"} to zero.

Secondly, we will define our custom color classes. Open \texttt{service.cpp}
again and add a enumeration definition for your language right after
the language information. You should start the
enumeration with the last default color class and still use the default color
classes where appropiate. The Haskell sample defines:
\begin{alltt}
enum HaskellColorClass 
\char123{}
  ClassPragma      = ClassDefaultLast + 1,

  ClassOperator,
  ClassConsId,
  ClassConsOp,

  ClassSpecialOp,
  ClassReservedOp,

  ClassGreencard,

  ClassLiterateText,
  ClassLiterateStart,

  //some short hands
  ClassSpecial     = ClassSpecialOp,	
  ClassReservedId  = ClassKeyword,
  ClassSpecialId   = ClassKeyword,
  ClassExtensionId = ClassKeyword,	
  ClassCharacter   = ClassString,	
  ClassSelector	   = ClassIdentifier,
  ClassImplicit	   = ClassIdentifier
\char125{};
\end{alltt}

In the token definition array, we can now use the new custom color classes. Here is an
excerpt from the Haskell sample:
\begin{alltt}
override const TokenInfo* Service::getTokenInfo() const
\char123{}
  static TokenInfo tokenInfoTable[] =
  \char123{}
    ...      
    \char123{} VAROP    , ClassOperator	, "operator (\%s)",    CharIdentifier \char125{},
    \char123{} VARID    , ClassIdentifier, "identifier (\%s)",  CharIdentifier \char125{},
    \char123{} CONOP    , ClassConsOp	, "constructor (\%s)", CharIdentifier \char125{},
    \char123{} CONID    , ClassConsId	, "constructor (\%s)", CharIdentifier \char125{},
    \char123{} QVAROP   , ClassOperator	, "operator (\%s)",    CharIdentifier \char125{},
    \char123{} QVARID   , ClassIdentifier, "identifier (\%s)",  CharIdentifier \char125{},
    \char123{} QCONOP   , ClassConsOp	, "constructor (\%s)", CharIdentifier \char125{},
    \char123{} QCONID   , ClassConsId	, "constructor (\%s)", CharIdentifier \char125{},
    \char123{} RECSELID , ClassSelector	, "selector",         CharIdentifier \char125{}, 
    \char123{} IPVARID  , ClassImplicit	, "implicit parameter", CharIdentifier \char125{},
    ...
  \char125{};
  return tokenInfoTable;
\char125{}
\end{alltt}

The last step consists of adding default color information for the new custom colors. 
The method \texttt{Service::getColorInfo} returns a pointer to a static array of
color information \texttt{ColorInfo}. Each \texttt{ColorInfo} consists of the
color class, its name and a CSS style specification of its appearance. 
The \texttt{getColorInfo} method should be added right after the custom color
class enumeration. The {\itshape{}Haskell} sample has the following implementation:
{\small
\begin{alltt}
override const ColorInfo* Service::getColorInfo() const
\char123{}
 static ColorInfo colorInfoTable[] =
 \char123{}
  \char123{} ClassKeyword,   "Keyword",    "color: blue; font-weight: bold"  \char125{},
  \char123{} ClassComment,   "Comment",    "background-color: lightgray"     \char125{},
  \char123{} ClassIdentifier,"Identifier", ""                                \char125{},
  \char123{} ClassString,    "String",     "color: purple"                   \char125{},
  \char123{} ClassNumber,    "Number",     ""                                \char125{},
      
  \char123{} ClassPragma,    "Pragma",                   "color: aqua"       \char125{},
  \char123{} ClassOperator,  "Haskell Operator",         ""                  \char125{},
  \char123{} ClassConsId,    "Haskell Constructor",      "font-weight: bold" \char125{},
  \char123{} ClassConsOp,    "Haskell Cons Operator",    "font-weight: bold" \char125{},
  \char123{} ClassSpecialOp, "Haskell Special Operator", ""                  \char125{},
  \char123{} ClassReservedOp,"Haskell Reserved Operator","color: purple"     \char125{},
  \char123{} ClassGreencard, "Greencard directive",      
                 "color: darkblue; text-decoration: line-through"   \char125{},

  \char123{} ClassLiterateText, "Haskell Literate Text", "color: darkgreen"  \char125{},    
  \char123{} ClassLiterateStart,"Haskell Literate Start","color: darkgray"   \char125{},

  \char123{} ColorClassEnd,     "<unknown>",       ""                  \char125{}
 \char125{};

 return colorInfoTable;
\char125{}
\end{alltt}
}

And we should add this method to the class declaration in \texttt{service.h}:
\begin{alltt}
class Service : public CommentService
\char123{}
protected:
  override const ColorInfo* getColorInfo() const;
  override const TokenInfo* getTokenInfo() const;
\char125{};
\end{alltt}

More information about CSS style specifications can be found on
\url{http://www.w3c.org} Allthough it is possible to specify any legal CSS
style, only a small subset is recognised by Visual Studio. Here is a list
of recognised properties and values:




%<![HTML[
%<table border=1 cellpadding=3>
%<tr><td><strong>property</strong></td><td><strong>values</strong></td>
%</tr>
%<tr valign=top><td><strong>font-weight</strong></td>
%<td>
%default<br>
%normal<br>
%bold<br>
%bolder<br>
%lighter<br>
%100<br>
%200<br>
%300<br>
%400<br>
%500<br>
%600<br>
%700<br>
%800<br>
%900
%</td>
%</tr>
%
%<tr valign=top><td><strong>text-decoration</strong></td>
%<td>
%default<br>
%normal<br>
%line-through<br>
%strike-through<br>
%</td></tr>
%
%<tr valign=top><td><strong>color</strong><br><strong>background-color</strong></td>
%<td>
%default<br>
%normal<br>
%auto<br>
%selected<br>
%inactive<br>
%plaintext<br>
%margin<br>
%black<br>
%white<br>
%maroon<br>
%darkgreen<br>
%brown<br>
%darkblue<br>
%purple<br>
%aquamarine<br>
%lightgray<br>
%darkgray<br>
%red<br>
%green<br>
%yellow<br>
%blue<br>
%magenta<br>
%cyan<br>
%lime<br>
%aqua<br>
%silver<br>
%olive<br>
%teal<br>
%navy
%</td></tr>
%<tr valign=top><td><strong>text-kind</strong></td>
%<td>
%default<br>
%normal<br>
%humantext<br>
%programtext
%</td></tr>
%</table>
%]]>



\begin{tabular}{|l|l|}
\hline
\textbf{color}
& default \\
\textbf{background-color}
& normal \\
& auto \\
& selected \\
& inactive \\
& plaintext \\
& margin \\
& black \\
& white \\
& maroon \\
& darkgreen \\
& brown \\
& darkblue \\
& purple \\
& aquamarine \\
& lightgray \\
& darkgray \\
& red \\
& green \\
& yellow \\
& blue \\
& magenta \\
& cyan \\
\hline
& lime \\
& aqua \\
& silver \\
& olive \\
& teal \\
& navy \\
\hline
\end{tabular}


\begin{tabular}{|l|l|}
\hline
\textbf{text-decoration}
& default \\
& normal \\
& line-through \\
& strike-through \\
\hline 
\end{tabular}


\begin{tabular}{|l|l|}
\hline
\textbf{text-kind}
& default \\
& normal \\
& humantext \\
& programtext \\
\hline
\end{tabular}

\begin{tabular}{|l|l|}
\hline 
\textbf{font-weight} 
& default \\
& normal \\
& bold  \\
& bolder \\
& lighter \\
& 100 \\
& 200 \\
& 300 \\
& 400 \\
& 500 \\
& 600 \\
& 700 \\
& 800 \\
& 900 \\
\hline
\end{tabular}


\hypertarget{Block commenting}{\section{Block commenting}}\label{Block commenting}

Using \texttt{ctrl-k, ctrl-c} or \texttt{ctrl-k, ctrl-u} in Visual Studio 7,
you can comment or uncomment selected text. This feature is not
supported in the Visual J++ 6.0 environment. To enable block commenting you
should implement the \texttt{Service::getCommentFormat} method.

First open \texttt{service.h} and add a declaration of the \texttt{getCommentFormat}
method:
\begin{alltt}
class Service : public CommentService
\char123{}
protected:
  override const CommentFormat* getCommentFormat() const;
  override const ColorInfo*     getColorInfo()     const;
  override const TokenInfo*     getTokenInfo()     const;
\char125{};
\end{alltt}

The method returns
a pointer to a \texttt{CommentFormat} structure which is defined in 
\texttt{client\char92{}common\char92{}stdservice.h}:
\begin{alltt}
struct CommentFormat
\char123{}
  const char* lineStart;
  const char* blockStart;
  const char* blockEnd;
  bool        useLineComments; 
\char125{};
\end{alltt}

The \texttt{lineStart} string should be the character sequence that
starts line comments, or \texttt{NULL} if not supported. \texttt{blockStart}
and \texttt{blockEnd} are the block comment start and end characters.
The \texttt{useLineComments} flag should be true if line comment characters
are preferred to block comments. This is used when the user selects a
block of lines; these can be commented using either line- or block comments.

Open \texttt{service.cpp} and implement this method. The MyC implementation is as follows:
\begin{alltt}
override const CommentFormat* Service::getCommentFormat() const
\char123{}
  static CommentFormat commentFormat =  \char123{} "//", "/*", "*/", true  \char125{};
  return \&commentFormat;
\char125{}
\end{alltt}

\hypertarget{Automatic brace matching}{\section{Automatic brace matching}}\label{Automatic brace matching}

When automatic brace matching is enabled, the language service will
highlight matching braces whenever the user types them\footnote{This
feature is not supported in the Visual J++ 6.0 environment}. The code in
the babels service also handles brace `jumping' when the user
types \texttt{ctrl-]} in Visual Studio.

\hypertarget{Properties}{\subsection{Properties}}\label{Properties}
First of all, the property \texttt{MatchBraces} should be set to one to
enable brace matching. When the property \texttt{ShowMatchingBrace} is
set to one, Visual Studio will show the line of the matching brace in 
the status bar; useful when the matching brace is not visible in the view.
The property \texttt{MatchBracesAtCaret} determines if the environment
will highlight matching braces even when just positioning the caret
on a brace, normally it will only highlight when a brace is typed. Be 
careful before enabling this feature since it might make the user experience
of the environment more distracting.

\hypertarget{Adding brace match rules}{\subsection{Adding brace match rules}}\label{Adding brace match rules}
Next, we are going to add grammar rules to \texttt{parser.y}
that specify when matching 
braces are found. The standard service \texttt{g\_service} provides the
following methods to signal matching braces:
\begin{alltt}
void matchPair  ( in const Location\& loc1, 
                  in const Location\& loc2 ) const;
void matchTriple( in const Location\& loc1, 
                  in const Location\& loc2, 
                  in const Location\& loc3 ) const;
\end{alltt}

Simple call one of these with the locations of the braces whenever
a matching pair (or triple!) is found. Here is for example
the code to match parenthesised expressions from the MyC sample:
\begin{alltt}
ParenExpr
    : '(' Expr ')'   \char123{} g\_service->matchPair(\$1,\$3); \char125{}
    | '(' Expr error 
          \char123{} g\_service->expectError( "unmatched parenthesis", ")" ); \char125{}
    ;
\end{alltt}

Off course, you should match braces whenever appropiate, also for example
when parsing a parameter list:
\begin{alltt}
ParenParams
    :  '(' ')'            \char123{} g\_service->matchPair(\$1,\$2); \char125{}
    |  '(' Params1 ')'    \char123{} g\_service->matchPair(\$1,\$3); \char125{}
    |  '(' Params1 error  \char123{} g\_service->expectError( "unmatched parenthesis", ")" ); \char125{}
    |  '(' error ')'      \char123{} g\_service->matchPair(\$1,\$3); 
                            g\_service->syntaxError( "parameters", \&\$2 ); \char125{}
\end{alltt}

Arbitrary ranges of text can be braces. Even triples can be matched as the
following code for recognising \texttt{if-else-then} statements from 
the Forth sample shows:
\begin{alltt}
IfBlock
  : KWIF IfWordList KWTHEN \char123{} g\_service->matchPair(\$1,\$3); \char125{}
  | KWIF IfWordList KWELSE IfWordList KWTHEN  
                           \char123{} g\_service->matchTriple(\$1,\$3,\$5); \char125{}
  | KWIF IfWordList error  \char123{} g\_service->expectError( "IFBlock", "THEN" ); \char125{}
  | KWIF IfWordList KWELSE IfWordList error   
                           \char123{} g\_service->matchPair(\$1,\$3);
                             g\_service->expectError( "IFBlock", "THEN" ); \char125{}
\end{alltt}

\hypertarget{Adding triggers}{\subsection{Adding triggers}}\label{Adding triggers}
Due to speed reasons, it is not possible to run the parser on each 
editor command. In order to recognise matching braces the babel service
needs to know when a parse might be needed, ie. when a token is entered
that looks like a brace token. The babel service implements this using
{\itshape{}triggers}. Triggers are specified in the \texttt{Service::getTokenInfo}
method in the \texttt{service.cpp} file. 

Open this file and add the trigger \texttt{TriggerMatchBraces} to each
token which is (or might be) a brace for your language.
Some relevant entries in the MyC sample are:
\begin{alltt}
override const TokenInfo* Service::getTokenInfo() const
\char123{}
  static TokenInfo tokenInfoTable[] =
  \char123{}
    //TODO: add your own token information here
    \char123{} IDENTIFIER, ClassIdentifier,  "identifier '\%s'",  CharIdentifier \char125{},
    ...
    \char123{} KWVOID, ClassKeyword, "void", CharKeyword   \char125{},
    \char123{} '(',    ClassText,  "'('"   , CharDelimiter, TriggerMatchBraces \char125{}, 
    \char123{} ')',    ClassText,  "')'"   , CharDelimiter, TriggerMatchBraces \char125{},
    \char123{} '\char123{}',    ClassText,  "'\char123{}'"   , CharDelimiter, TriggerMatchBraces \char125{},
    \char123{} '\char125{}',    ClassText,  "'\char125{}'"   , CharDelimiter, TriggerMatchBraces \char125{},
    \char123{} '.',    ClassText,  "'.'"   , CharDelimiter \char125{},
    ...
    //always end with the 'TokenEnd' token.
    \char123{} TokenEnd,     ClassText,      "<unknown>"   \char125{}
  \char125{};

  return tokenInfoTable;
\char125{};
\end{alltt}

The Forth language service will for example add the trigger too 
for the \texttt{KWIF}, \texttt{KWTHEN} and \texttt{KWELSE} tokens.


\hypertarget{Statement completion}{\section{Statement completion}}\label{Statement completion}

{\itshape{}warning: these interfaces might still change}



Statement completion shows up in many ways in Visual Studio. The traditional
statement completion is invoked whenever you type 
\texttt{ctrl-<left arrow>} or \texttt{ctrl-<space>}. The environment
will complete the current word if possible or show a list possible candidates
if there is ambiguity. The key combination \texttt{ctrl-j} always shows the
list of identifiers in scope. 

\hypertarget{Properties}{\subsection{Properties}}\label{Properties}
The property \texttt{ShowCompletion} determines if the options
{\itshape{}show member list} and {\itshape{}show parameter list} are visible
in the Visual Studio environment. This property is enabled by default.
The property \texttt{SortMemberList} determines if the identifier
list is alphabetized. This is enabled by default.

\hypertarget{Giving scope information}{\subsection{Giving scope information}}\label{Giving scope information}
In order to do statement completion, information about the declarations
and scopes in the source file has to be known to the system.
The most general and powerful method to
give scope information is by implementing the \texttt{IScope} interface.

However, the babel package provides a default implementation that can
deal with almost any language is satisfying manner. This document only
describes this approach.

The standard service \texttt{g\_service} has two methods to give scope
and declaration information:
\begin{alltt}
void addScope( in const Location\& start, in const Location\& end,
               in ScopeKind       kind, 
               in ScopeAccess     access, 
               in ScopeStorage    storage,
               in const Location\& name, 
               in const Location\& descStart, 
               in const Location\& descEnd,
               in const Location* type = NULL,             
               in bool            merge = false 
              );
void addScopeText( in const Location\& start, in const Location\& end,
                   in ScopeKind     kind, 
                   in ScopeAccess   access, 
                   in ScopeStorage  storage,
                   in const char*   name    = NULL, 
                   in const char*   description = NULL,
                   in const char*   type    = NULL,
                   in const char*   display = NULL, 
                   in bool          merge = false );
\end{alltt}

The first method takes raw locations for descriptions while the second
one takes arbitrary text. The following functions can be used to combine
text locations or to get a pointer to the text at a location:
\begin{alltt}
Location    range( in const Location\& loc1, 
                   in const Location\& loc2 ) const;
const char* tokenText( in const Location* = NULL, 
                       in const Location* = NULL ) const;   
\end{alltt}

(The pointer returned by \texttt{tokenText} is in \texttt{static} memory).
For each declaration and scope, the method \texttt{addScope} (or
\texttt{addScopeText}) is called with the start and ending location. These
locations give the lexical scope of this declarations. If a declaration
doesn't have a lexical scope, like  a variable declaration, the locations
might simply coincide. The \texttt{kind} determines if this declaration
has parameters or if it can contain other scopes; it contains no other
semantic information. The kinds are declared in \texttt{babelservice.idl}:
\begin{alltt}
enum ScopeKind
\char123{}
  ScopeUnknown = 0,
  ScopeModule,
  ScopeClass,
  ScopeInterface,
  ScopeUnion,

  ScopeProcedure,
  ScopeVariable,

  ScopeBlock
\char125{};
\end{alltt}

The \texttt{access} determines the visibility of the declaration and
its members:
\begin{alltt}
enum ScopeAccess
\char123{}
  AccessPrivate = 0,
  AccessProtected,
  AccessPublic
\char125{};
\end{alltt}

The \texttt{storage} determines, together with the kind, the icon that
is displayed for this declaration and it is used to find member information.
\begin{alltt}
enum ScopeStorage
\char123{}
  StorageConstant = 0,
  StorageStatic,
  StorageMember,
  StorageParameter,
  StorageResult,
  StorageLocal,
  StorageType,
  StorageOther
\char125{};
\end{alltt}

The \texttt{name} contains the name of this declaration. The 
\texttt{descStart} and \texttt{descEnd} parameters contain the start
and end location of the declaration. The text between those location
is the description of this declaration. The \texttt{type} location 
contains the type of this declaration. The \texttt{merge} flag should be
set to \texttt{true} if the declaration is specified in multiple places
and the information should be merged with the previous entries.

The rules in \texttt{parser.y} should now be extended with calls
to \texttt{addScope}. Let's look at some of the rules in the MyC sample:
\begin{alltt}
Declaration
  : Type IDENTIFIER ParenParams Block                                                                       
    \char123{} g\_service->addScope( \$1, \$4, 
                           ScopeProcedure, AccessPublic, StorageType,
                           \$2, \$1, \$3, \&\$1 ); \char125{} 
  ;
\end{alltt}

This rule describes a declaration of a public procedure in MyC.
The scope of the method extends from the start of \texttt{Type} to
the end of \texttt{Block}, ie. \texttt{\$1} to \texttt{\$4}.
The name is given by \texttt{IDENTIFIER}, ie. \texttt{\$2}.
The description that should be shown in the quick info, is just
the declaration without the block of code, ie. \texttt{\$1} to \texttt{\$3}.



By default, the location of a production is the same as that of its first
token; each rule contains the implicit assignment: \texttt{\$\$ = \$0;}.
Since we want to make sure that the location of \texttt{Block} entails the
entire block (and not just the first token \texttt{\char123{}}) we need to add some
rules to the \texttt{Block} too:
\begin{alltt}
Block
  : '\char123{}' '\char125{}'   
      \char123{} \$\$ = g\_service->range(\$1,\$2);
        g\_service->matchPair(\$1,\$2); 
        g\_service->addScopeText( \$1, \$2, 
                                 ScopeBlock, AccessPublic, StorageOther); 
      \char125{}
  | '\char123{}' BlockContent1 '\char125{}'
      \char123{} \$\$ = g\_service->range(\$1,\$3);
        g\_service->matchPair(\$1,\$3);
        g\_service->addScopeText( \$1, \$3, 
                                 ScopeBlock, AccessPublic, StorageOther);  
      \char125{}
  ;
\end{alltt}

The first action of the rules is to assign the whole range from begin 
to end to the location of the \texttt{Block} production. In this particular
case it would even be good enough to do \texttt{\$\$ = \$2} and \texttt{\$\$ = \$3}
respectively but using \texttt{range} is good coding style.



A block itself imposes a lexical scope allthough it has no name or description.
Here the method \texttt{addScopeText} is called with the \texttt{ScopeBlock}
kind and no text for the name, description or type.

\hypertarget{Member completion}{\section{Member completion}}\label{Member completion}

Most languages have some kind of modules, classes or structures
that have members. These members are normally selected using some kind of
operator, like the \texttt{'.'} in Java or the \texttt{'->'} in C++.
The grammar rules need to be extended let statement completion work 
correctly when members are involved. The service has two methods that
specify how a member is selected:
\begin{alltt}
void startSelect( in const Location\& name) const;
void addSelect( in const Location\& selector, 
                in const Location\& name) const;
\end{alltt}

The grammar should call \texttt{startSelect} whenever a identifier is
used in an expression. The method \texttt{addSelect} should be called for
every identifier that is selected. For example, a Java grammar should
contain the following rules for {\itshape{}qualified} names:
\begin{alltt}
QualifiedName
  : IDENTIFIER                    \char123{} g\_service->startSelect(\$1); \char125{}
  | QualifiedName '.' IDENTIFIER  \char123{} g\_service->addSelect(\$2,\$3); \char125{} 
  | QualifiedName '.' error       \char123{} g\_service->addSelect(\$2,\$2); 
                                    g\_service->expectError( "qualified name","identifier" ); \char125{}
\end{alltt}

The last production with the \texttt{error} is important in practice since
this case will occur quite often when the user is typing its source. In
the case where the name is completely lacking, the location of the selector
should be passed (\texttt{\$2} in the example).

\hypertarget{Adding member list triggers}{\subsection{Adding member list triggers}}\label{Adding member list triggers}
When typing a selector token, the member selection list should automatically 
pop up (unless disabled in the option settings of Visual Studio).
Again this is implemented by triggers. The trigger \texttt{TriggerMemberSelect}
triggers a member completion list.
Here is a typical example of such entry in the \texttt{getTokenInfo} method:
\begin{alltt}
 ...
 \char123{} '.',  ClassText,  "'.'", CharDelimiter, TriggerMemberSelect \char125{},   
 ...
\end{alltt}

\hypertarget{Quick Info}{\section{Quick Info}}\label{Quick Info}
Quick info is a feature that will pop up a small text box that describes
an identifier when the user hovers the mouse above it. The property
\texttt{QuickInfo} should be set to zero to disable quick info. It is
enabled by default. Quick info will work when you have member completion 
implemented. Even when the language doesn't have member selection, you
should still call \texttt{startSelect} at every identifier expression.

\hypertarget{Method tips}{\section{Method tips}}\label{Method tips}
Method tips are windows with parameter information about the procedure
or method that is called. A method tip normally pops up whenever the
user types a parameter list token, for example an open parenthesis.
It can also be invoked manually by typing \texttt{ctrl-shift-<space>}
\footnote{\texttt{ctrl-shift-I} in VJ++ 6.0}. 
Method tips are automatically supported by Babel when parameter information
is added to the parser rules.

\hypertarget{Adding parameter information}{\subsection{Adding parameter information}}\label{Adding parameter information}
The global service has three methods that support parameter information.
\begin{alltt}
void startParameters( in const Location\& ) const;
void parameter(in const Location\& ) const;
void endParameters(in const Location\& ) const;
\end{alltt}

The method \texttt{startParameters} is called whenever a parameter list
starts with the location of the parameter list start token. \texttt{parameter}
is called for every parameter seperator found and \texttt{endParameters}
when the parameter list is terminated. Here is a typical example from MyC:
\begin{alltt}
ParenArguments
    : StartArg EndArg            \char123{} g\_service->matchPair(\$1,\$2); \char125{} 
    | StartArg Arguments1 EndArg \char123{} g\_service->matchPair(\$1,\$3); \char125{}
    | StartArg Arguments1 error  
              \char123{} g\_service->endParameters(@3);
                g\_service->expectError( "unmatched parenthesis", ")" ); \char125{}
    ;

StartArg
    : '('     \char123{} g\_service->startParameters(\$1); \char125{}
    ;

EndArg
    : ')'     \char123{} g\_service->endParameters(\$1); \char125{}
    ;    

Arguments1
    : Expr ','  \char123{} g\_service->parameter(\$2); \char125{}  Arguments1
    | Expr
    ;
\end{alltt}

In MyC, parenthesis enclose a parameter list. The example is complicated
by the need to call the different methods exactly at the moment one of 
those tokens is recognised. In the rule \texttt{Arguments1} a mid-rule
action is needed in order to call \texttt{parameter} at the right moment.
Suppose that rule would have been written like this:
\begin{alltt}
Arguments1
    : Expr ','  Arguments1  \char123{} g\_service->parameter(\$2); \char125{}  
    | Expr
    ;
\end{alltt}
In this case \texttt{parameter} will called {\itshape{}after} the 
other arguments in the production \texttt{Arguments1} are recognised.
Now all \texttt{parameter}'s calls are done in a reverse order! 

The same kind of problem occura when calling \texttt{startParameter} and
\texttt{endParameter}, that is why the \texttt{StartArg} and \texttt{EndArg}
productions are explicitly added.

\hypertarget{Adding parameter triggers}{\subsection{Adding parameter triggers}}\label{Adding parameter triggers}
Allthough methods tips can now be explicitly invoked, they are much more
useful when they pop up automatically when the user opens a parameter list.
The triggers \texttt{TriggerParamStart}, \texttt{TriggerParamEnd} and
\texttt{TriggerParamNext} trigger a method tip.
Simply add the relevant entries to the token info table. The
entries for MyC look like:
\begin{alltt}
   ...
   \char123{} '(', ClassText,  "'('", 
          CharDelimiter, TriggerParamStart | TriggerMatchBraces \char125{}, 
   \char123{} ')', ClassText,  "')'", 
          CharDelimiter, TriggerParamEnd | TriggerMatchBraces   \char125{},
   \char123{} ',', ClassText,  "','", CharDelimiter, TriggerParamNext    \char125{},     
   ...
\end{alltt}

Not all languages distinguish the parameter list open token from
the parameter seperator or closing token. It is always legal to combine
triggers. Allthough triggers will trigger a parser action in order
to show or dismiss a method tip, the distinction in triggers is only used
to optimize some common cases. The parser always has the final decision
in showing the appropiate information.

For example, in the language Haskell, white space can open a 
parameter list but also seperates parameters. At the same time there is
no specific token for closing a parameter list but a whole set of tokens
can close a parameter list. All those tokens should get the \texttt{TriggerParamEnd}
trigger while white space gets two triggers:
\begin{alltt}
  ...
  \char123{} LEX\_WHITE, ClassText, "white space", 
               CharWhiteSpace, TriggerParamStart | TriggerParamNext \char125{},
  ...
\end{alltt}

In the Haskell case, there is a more serious problem. Since white space
now triggers the parameter info, the parser is called on each white space
character typed which will degrade the responsiveness of the editor.
The solution is to recognise `special' white space that might trigger
a method tip in the lexer. In Haskell, that is the first white space after
a valid identifier. The following rules in \texttt{lexer.lex} recognise
this case:
\begin{alltt}
<init>\char123{}Identifier\char125{}    \char123{} g\_service->setLexState(afterid); 
                        return IDENTIFIER; \char125{}

<afterid>\char123{}White\char125{}      \char123{} g\_service->setLexState(init);
                        return LEX\_SPECIAL\_WHITE; \char125{}
<afterid>.            \char123{} g\_service->setLexState(init);
                        yyless(0); \char125{}
\end{alltt}


\hypertarget{IDL specification}{\section{IDL specification}}\label{IDL specification}
{\small
\begin{alltt}
/*---------------------------------------------------------
BabelServiceLib:
  Interface specification for the Babel language services.  
  
  Midl will convert this file to 
  - a type library [babelservice.tlb]
  - a header file [babelservice.h]
  - a GUID definition file [babelservice\_i.c].
---------------------------------------------------------*/  
[
  uuid(4E71BA42-5C36-11d3-B3FD-006008D1BF8C),
  version(1.0),
  helpstring("BabelService 1.0 Type Library")
]
library BabelServiceLib
\char123{}
  importlib("stdole32.tlb");
  importlib("stdole2.tlb");

  /*---------------------------------------------------------
  Severity:
    The severity of an error. Used for example to select
    an appropiate bitmap in the task window.
  ---------------------------------------------------------*/
  enum Severity 
  \char123{}
    SevHint,
    SevWarning,
    SevError,
    SevFatal
  \char125{};

  /*---------------------------------------------------------
  ParseReason:
    The reason why the parser is called.
  ---------------------------------------------------------*/
  enum ParseReason
  \char123{}
	ReasonColorize,
    ReasonCheck,
    ReasonMemberSelect,
    ReasonCompleteWord,
    ReasonQuickInfo,
    ReasonMethodTip,
    ReasonMatchBraces
  \char125{};

  /*---------------------------------------------------------
  CharClass:
    The class of the current character. Used for example
    to determine the word extent and search ranges.
  ---------------------------------------------------------*/
  enum CharClass
  \char123{}
    CharDefault     = 0,
    CharText        = CharDefault,
    CharKeyword,
    CharIdentifier,
    CharString,
    CharLiteral,

    CharOperator,
    CharDelimiter,
    CharEmbedded,

    CharWhiteSpace,
    CharLineComment,
    CharComment,

    CharDefaultLast = CharComment
  \char125{};

  /*---------------------------------------------------------
  TriggerClass:
    If a character has (a) trigger(s) associated with it, it may
    fire one of the following actions when it is entered by a user:
    [TriggerMemberSelect]  a member selection tip window
    [TriggerMethodTip]     a method tip window
    [TriggerMatchBraces]   highlight matching braces
  ---------------------------------------------------------*/
  typedef long TriggerClass;

  enum DefaultTriggerClass
  \char123{}
    TriggerNone             = 0,
    TriggerMemberSelect     = 0x01,
    TriggerMatchBraces      = 0x02,

    TriggerMethodTip        = 0xf0,
    TriggerParamStart       = 0x10,
    TriggerParamNext        = 0x20,
    TriggerParamEnd         = 0x40,
    TriggerParam            = 0x80
  \char125{};

  /*---------------------------------------------------------
  ColorClass:
    Determines the color of a character.
  ---------------------------------------------------------*/
  typedef long ColorClass;

  /*---------------------------------------------------------
  DefaultColorClass:
    These default color classes are always available. If a 
    client doesn't need custom color classes, it can set
    the registry entry "RequestStockColors" to 1 and use
    these defaults.
  ---------------------------------------------------------*/
  enum DefaultColorClass
  \char123{}
    ClassDefault    = 0,
    ClassText       = ClassDefault,
    ClassKeyword,
    ClassComment,
    ClassIdentifier,
    ClassString,
    ClassNumber,
    ClassDefaultLast = ClassNumber,
  \char125{};

  /*---------------------------------------------------------
  ScopeKind:
    The kind of scope that is opened. These kinds have no
    semantic significance. They are only used to determine
    the bitmaps used when displaying scope information
  ---------------------------------------------------------*/
  enum ScopeKind
  \char123{}
    ScopeUnknown = 0,
    ScopeModule,
    ScopeClass,
    ScopeInterface,
    ScopeUnion,

    ScopeProcedure,
    ScopeVariable,

    ScopeBlock,
  \char125{};

  enum ScopeAccess
  \char123{}
    AccessPrivate = 0,
    AccessProtected,
    AccessPublic
  \char125{};

  enum ScopeStorage
  \char123{}
    StorageConstant = 0,
    StorageStatic,
    StorageMember,
    StorageParameter,
    StorageResult,
    StorageLocal,
    StorageType,
    StorageOther
  \char125{};

         
  /*---------------------------------------------------------
  IColorSink:
    This interface is passed as an argument to the [ColorLine]
    method. Its main use is to set the color class of each
    character but it is also be used to pass the character class
    and trigger information.
  ---------------------------------------------------------*/
  [ object
  , uuid(4E71BA43-5C36-11d3-B3FD-006008D1BF8C)
  , helpstring("IColorSink Interface")
  , oleautomation
  , dual
  ]
  interface IColorSink : IDispatch
  \char123{}       
    /*---------------------------------------------------------
    Colorize:   
      Set the colorclasses for a range of characters in a line.
      The [startIdx] and [endIdx]  arguments specify the character 
      index in the line (in constrast to columns, this means that 
      '\char92{}t' characters are treated as single character). 
      The indices are 0 based, ie. the first character has index 0.
    ---------------------------------------------------------*/
    [helpstring("method Colorize")] 
    HRESULT Colorize( [in] long startIdx,
                      [in] long endIdx,
                      [in] ColorClass colorClass,
                      [in] enum CharClass charClass,
                      [in] TriggerClass trigger );
  \char125{};


  /*---------------------------------------------------------
  IParseSink:
    This interface is passed as an argument to the [ParseSource]
    function. It is used to return information gathered
    during parsing to the VS environment. 
    It is mainly used to send (warning/error) messages.
  ---------------------------------------------------------*/
  [ object
  , uuid(4E71BA40-5C36-11d3-B3FD-006008D1BF8C)
  , helpstring("IParseSink Interface")
  , oleautomation
  , dual
  ]
  interface IParseSink : IDispatch
  \char123{}
    /*---------------------------------------------------------
    ErrorMessage:
      Send a [message] to the VS enviroment. The kind of message
      is specified through the severity [sev]. Note that
      lines and indices are all 0 based.
    ---------------------------------------------------------*/
    [helpstring("method ErrorMessage")] 
    HRESULT ErrorMessage( [in] long startLine, [in] long endLine, 
                          [in] long startIdx,  [in] long endIdx,
                          [in] enum Severity sev, [in] BSTR message );

    /*---------------------------------------------------------
    MatchPair:
      Whenever a matching pair is parsed, this method is called
      with the text span of both the left and right item. The
      information is used to when a user types "ctrl-]" in VS
      to find a matching brace and when auto-highlight matching
      braces is enabled.
    ---------------------------------------------------------*/    
    HRESULT MatchPair( [in] long startLine1, [in] long startIdx1,
                       [in] long endLine1,   [in] long endIdx1, 
                       [in] long startLine2, [in] long startIdx2,
                       [in] long endLine2,   [in] long endIdx2 );

    HRESULT MatchTriple( [in] long startLine1, [in] long startIdx1,
                         [in] long endLine1,   [in] long endIdx1, 
                         [in] long startLine2, [in] long startIdx2,
                         [in] long endLine2,   [in] long endIdx2,
                         [in] long startLine3, [in] long startIdx3,
                         [in] long endLine3,   [in] long endIdx3 );

    /*---------------------------------------------------------
    Member Selection:
      [StartSelect] is called for each name that is parsed ("Stack").
      [AddSelect] is called for each qualification with both
      the text span of the selector (".")  and the text span 
      of the name ("Push").
    ---------------------------------------------------------*/    
    HRESULT StartSelect( [in] long line, [in] long startIdx, [in] long endIdx );
    HRESULT AddSelect( [in] long lineSelector,
                       [in] long startIdxSelector, [in] long endIdxSelector,
                       [in] long line, [in] long startIdx, [in] long endIdx );

    /*---------------------------------------------------------
    Method Info:
      [StartParameters] is called when the parameters of a method
                        are started "(".
      [Parameter] is called on the start of a new parameter ",".
      [EndParameter] is called on the end of the paramters ")".
    ---------------------------------------------------------*/    
    HRESULT StartParameters( [in] long line, [in] long idx );
    HRESULT Parameter      ( [in] long line, [in] long idx );
    HRESULT EndParameters  ( [in] long line, [in] long idx );


    /*---------------------------------------------------------
    Scope Builder:
      The babel package provides a default implementation
      of the [IScope] interface which can be used through the
      following methods. If you want to provide your own
      [IScope] implementation, don't call these methods and
      return the scope trough the [scope] parameter of the
      [IBabelService::ParseSource] method.
    ---------------------------------------------------------*/            
    HRESULT AddScope( [in] long startLine, [in] long startIdx,
                      [in] long endLine,   [in] long endIdx,
                      [in] enum ScopeKind kind,
                      [in] enum ScopeAccess access,
                      [in] enum ScopeStorage storage,
                      [in] BSTR name,     [in] BSTR type,
                      [in] BSTR display,  [in] BSTR description,
                      [in] VARIANT\_BOOL merge
                    );                       
    HRESULT AddInclude( [in] enum ScopeAccess access, [in] BSTR name );
    HRESULT AddExtern( [in] BSTR external );
  \char125{};
  
  /*---------------------------------------------------------
  INames
  ---------------------------------------------------------*/    
  [ object
  , uuid(4E71BA45-5C36-11d3-B3FD-006008D1BF8C)
  , helpstring("INames Interface")
  , oleautomation
  , dual
  ]
  interface INames : IDispatch
  \char123{}
    HRESULT GetCount( [out,ref,retval] long* count );
    HRESULT GetName( [in] long index, [out,ref,retval] BSTR* name );
  \char125{};
  
  /*---------------------------------------------------------
  IDeclarations
  ---------------------------------------------------------*/    
  [ object
  , uuid(4E71BA46-5C36-11d3-B3FD-006008D1BF8C)
  , helpstring("IDeclarations Interface")
  , oleautomation
  , dual
  ]
  interface IDeclarations : IDispatch
  \char123{}
    HRESULT GetCount      ( [out,ref,retval] long* count );
    HRESULT GetName       ( [in] long index, [out,ref,retval] BSTR* name );
    HRESULT GetDescription( [in] long index, [out,ref,retval] BSTR* description );    
    HRESULT GetGlyph      ( [in] long index, [out,ref,retval] long* glyph );

    HRESULT GetBestMatch( [in] BSTR text, [out] long* index, [out] VARIANT\_BOOL* uniqueMatch );
  \char125{};
  

  /*---------------------------------------------------------
  IMethods
  ---------------------------------------------------------*/    
  [ object
  , uuid(4E71BA47-5C36-11d3-B3FD-006008D1BF8C)
  , helpstring("IMethods Interface")
  , oleautomation
  , dual
  ]
  interface IMethods : IDispatch
  \char123{}
    HRESULT GetName ( [out,ref,retval] BSTR* name  );
    HRESULT GetCount( [out,ref,retval] long* count );

    HRESULT GetDescription( [in] long index, [out,ref,retval] BSTR* description );    
    HRESULT GetType       ( [in] long index, [out,ref,retval] BSTR* type );
    
    HRESULT GetParameterCount( [in] long index, [out,ref,retval] long* count );
    HRESULT GetParameterInfo ( [in] long index, 
                               [in] long parameter, 
                               [out,unique] BSTR* name,
                               [out,unique] BSTR* display,
                               [out,unique] BSTR* description );
  \char125{};
  
  

  /*---------------------------------------------------------
  IScope:
  ---------------------------------------------------------*/    
  [ object
  , uuid(4E71BA44-5C36-11d3-B3FD-006008D1BF8C)
  , helpstring("IScope Interface")
  , oleautomation
  , dual
  ]
  interface IScope : IDispatch
  \char123{}
    HRESULT GetDeclarations( [in] long line, 
                             [in] long idx, 
                             [in,unique] INames* names,
                             [out,ref,retval] IDeclarations** decls );  

    HRESULT GetDataTipText ( [in] long line, 
                             [in] long idx, 
                             [in,unique] INames* names,
                             [out,ref,retval] BSTR* text );  

    HRESULT GetMethods     ( [in] long line,
                             [in] long idx,
                             [in,unique] INames* names,
                             [out,ref,retval] IMethods** imethods );
  \char125{};


  /*---------------------------------------------------------
  IBabelService:
    By implementing this interface, a client gets a full language 
    service automatically provided by Visual Studio.
  ---------------------------------------------------------*/
  [ object
  , uuid(4E71BA41-5C36-11d3-B3FD-006008D1BF8C)
  , helpstring("IBabelService Interface")
  , oleautomation
  , dual
  ]
  interface IBabelService : IDispatch
  \char123{}
    [helpstring( "ColorCount: return the number of custom color classes" )]
    HRESULT ColorCount( [out,ref,retval] ColorClass* count );

    /*---------------------------------------------------------
      GetColorInfo:
        retrieve color information for a specific color class.
        This method is only called if the registry entry
        [RequestStockColors] is set to 0.
      pre: 
        [index] < [ColorCount()]
      returns: 
        a [description] of the color class ("keyword")
        and a CSS [style] for this color class, for example
        "color: blue; font-weight: bold"
    ---------------------------------------------------------*/
    [helpstring("GetColorInfo: get color class information")]
    HRESULT GetColorInfo( [in]  ColorClass index, 
                          [out,ref] BSTR* description,
                          [out,ref] BSTR* style );

    /*---------------------------------------------------------
      ColorLine: 
        colorize a line [line]. Colors can be set through the
        [sink] that is passed too. The [state] can be used by
        the client in anyway it wants; it always corresponds to
        the state that was set at the end of the previous line.
        [state] is 0 if no state is known (start of the file).
      pre:          
        [sink] may be NULL in order to just retrieve the 
        new state at the end of the line
    ----------------------------------------------------------*/
    [helpstring("ColorLine: colorize a line")]    
    HRESULT ColorLine( [in] BSTR line, 
                       [in,unique] IColorSink* sink,
                       [in,out,ref] long* state 
                     );
    
    /*---------------------------------------------------------
      ParseSource:
        parses the program source in [text]. Messages containing
        warnings or errors can be sent to the [sink]. The reason
        for calling [ParseSource] is passed in [reason]; a client
        can always choose not to do anything.
    ---------------------------------------------------------*/
    [helpstring("ParseSource: parse the source")] 
    HRESULT ParseSource( [in] BSTR text, 
                         [in,unique] IParseSink* sink,
                         [in] enum ParseReason reason,
                         [out,unique,retval] IScope** scope 
                        );

    /*---------------------------------------------------------
      GetMethodFormat:
        Used for formatting method data tips.
        [typePrefixed] is [true] if the return type should be
        in front of the name: [long foo()] and [false]
        if it comes last: [foo() as long].
        [parStart],[parSep] and [parEnd] give the parameter start
        token, the seperator and end token respectively. 
        [typeStart] and [typeEnd] surround the return type.
        All parameters can be set to NULL in which case they
        default to a C-style formatting.
    ---------------------------------------------------------*/
    [helpstring("GetMethodFormat: return format tokens for method data tips")]
    HRESULT GetMethodFormat( [out,ref] BSTR* parStart,
                             [out,ref] BSTR* parSep, 
                             [out,ref] BSTR* parEnd,
                             [out,ref] BSTR* typeStart, 
                             [out,ref] BSTR* typeEnd,
                             [out,ref] VARIANT\_BOOL* typePrefixed );

    /*---------------------------------------------------------
      GetCommentFormat:
        Used for (un)commenting selected text.
        [lineStart] is the token used to start line comments (//)
        [blockStart] is the token used to start block comments
        [blockEnd] is the token used to end block comments
        if [useLineComments]  is true, selected text is commented
        using line comments, otherwise it is commented using
        block comments.
    ---------------------------------------------------------*/
    [helpstring("GetCommentFormat: return format tokens for comments")]
    HRESULT GetCommentFormat( [out,ref] BSTR* lineStart, 
                              [out,ref] BSTR* blockStart,
                              [out,ref] BSTR* blockEnd,
                              [out,ref] VARIANT\_BOOL* useLineComments );

    
    [helpstring("GetImageList: return HIMAGELIST for glyphs")]
    HRESULT GetImageList( [out] long* imageListHandle );
  \char125{};
\char125{};
\end{alltt}
}


\end{document}

